% =============================================================================
% APPENDIX BI: DOMAIN-CORPORATE-STRATEGY: Strategy Design Model v1.0
% =============================================================================
% Module: EBF_DOMAIN_CORPORATE_01
% Category: DOMAIN
% Language: English
% Version: 1.0 (January 2026)
% Status: Experimental
%
% Strategy Design Model v1.0 (SDM): A process-based framework for designing
% comprehensive corporate strategies using EBF Hybrid methodology (Theory + Practice).
% Models strategic transformation as 4-phase behavioral change journey with
% context-dependent complementarity effects.
%
% =============================================================================

%%%%%%%%%%%%%%%%%%%%%%%%%%%%%%%%%%%%%%%%%%%%%%%%%%%%%%%%%%%%%%%%%%%%%%%%%%%%%%%
% PART 1: FRONT MATTER
%%%%%%%%%%%%%%%%%%%%%%%%%%%%%%%%%%%%%%%%%%%%%%%%%%%%%%%%%%%%%%%%%%%%%%%%%%%%%%%

% =============================================================================
% CHAPTER-APPENDIX LINKAGE BOX
% =============================================================================

\begin{tcolorbox}[colback=purple!5!white, colframe=purple!75!black,
    title=Chapter Linkage: Corporate Strategy Design]

\textbf{Primary Application:} Chapter 14: Organization-Level Decision-Making
\begin{itemize}[nosep]
\item Section 14.3: Strategic Decision Frameworks $\leftrightarrow$ This Appendix Section 2.2
\item Section 14.5: Implementation and Change Management $\leftrightarrow$ This Appendix Section 3
\end{itemize}

\textbf{Secondary Applications:}
\begin{itemize}[nosep]
\item Chapter 5: Complementarities in Practice --- uses $\gamma$ matrix from Section 2.3
\item Chapter 11: Context-Dependent Decision-Making --- uses $\Psi$ dimensions from Section 2.1
\end{itemize}

\textbf{Reading Guide:}
\begin{itemize}[nosep]
\item For strategic intuition: Start with Section 1 (Fundamental Question)
\item For formal specification: Read Sections 2-3 (Theory and Model)
\item For implementation: Read Section 4 (Worked Examples)
\item For parameters: See BBB Appendix (Parameter Repository)
\end{itemize}

\end{tcolorbox}

% =============================================================================
% APPENDIX TITLE AND LABEL
% =============================================================================

\chapter{DOMAIN-CORPORATE-STRATEGY: Strategy Design Model v1.0}
\label{app:sdm}


\begin{tcolorbox}[colback=blue!5!white, colframe=blue!75!black,
    title=Quick Reference: Appendix UNMAPPED_COR]

\textbf{Appendix:} COR (DOMAIN)

\textbf{Core Question:} What are the key findings in this domain?

\textbf{Key Concepts:}
\begin{itemize}[nosep]
\item Framework integration
\item Parameter validation
\item Cross-references
\end{itemize}

\textbf{Cross-References:} See related appendices below
\end{tcolorbox}

% =============================================================================
% ABSTRACT
% =============================================================================

\begin{abstract}
This appendix presents the Strategy Design Model v1.0 (SDM), a process-based framework for designing comprehensive corporate strategies using the EBF hybrid methodology. SDM models strategic transformation as a four-phase behavioral change journey (Awareness-Readiness-Action-Embedding) over 3-5 years, integrating five context dimensions ($\Psi$), five strategic utility dimensions (Financial, Environmental, Practical, Strategic, Deliberative), and seven complementarity parameters ($\gamma$). The model uses Behavioral Change Journey (BCJ) mathematics to predict strategy adoption probability and outcome distributions across organizational change phases. SDM is designed for large organizations facing market transformation and serves as the bridge between EBF theory and organizational practice.
\end{abstract}

% =============================================================================
% QUICK REFERENCE BOX
% =============================================================================

\begin{tcolorbox}[colback=blue!5!white, colframe=blue!75!black,
    title=Quick Reference: Strategy Design Model v1.0]

\textbf{Core Question:} How do organizations transition from current strategy ($S_0$) to target strategy ($S_1$) over 3-5 years, accounting for complementarities between strategic dimensions and context factors?

\textbf{Key Formula (Strategic Utility):}
\begin{equation}
U_{\text{strategy}}(t) = \beta_0 + \sum_{i} \beta_i C_i(t) + \sum_{i<j} \gamma_{ij} C_i(t) C_j(t) + \Psi_t \cdot \lambda
\end{equation}

\textbf{Key Concepts:}
\begin{itemize}[nosep]
\item \textbf{BCJ (Behavioral Change Journey):} 4-phase model (Awareness, Readiness, Action, Embedding)
\item \textbf{Context Modulation ($\Psi_t$):} Time-varying environmental factors (market disruption, capital, leadership alignment)
\item \textbf{Complementarity ($\gamma_{ij}$):} Synergies (+) and trade-offs (-) between strategic dimensions
\item \textbf{Strategic Dimensions (C)}: FEPSDE framework adapted for corporate strategy
\end{itemize}

\textbf{Cross-References:} Appendix UNMAPPED_HOW (HOW: Complementarities), Appendix CTW (WHEN: Context), Appendix UNMAPPED_EST (Parameters), Appendix UNMAPPED_DSN (EEE Workflow used in SDM design)

\end{tcolorbox}

%%%%%%%%%%%%%%%%%%%%%%%%%%%%%%%%%%%%%%%%%%%%%%%%%%%%%%%%%%%%%%%%%%%%%%%%%%%%%%%
% PART 2: CORE CONTENT
%%%%%%%%%%%%%%%%%%%%%%%%%%%%%%%%%%%%%%%%%%%%%%%%%%%%%%%%%%%%%%%%%%%%%%%%%%%%%%%

%%%%%%%%%%%%%%%%%%%%%%%%%%%%%%%%%%%%%%%%%%%%%%%%%%%%%%%%%%%%%%%%%%%%%%%%%%%%%%%
% -----------------------------------------------------------------------------
% APPENDIX SCOPE BOX
% -----------------------------------------------------------------------------

\begin{tcolorbox}[
    colback=orange!5!white,
    colframe=orange!75!black,
    title={\textbf{Appendix Scope: Ziel / In-Scope / Out-of-Scope / Constraints}},
    fonttitle=\bfseries
]

\textbf{Ziel (Objective):}
\begin{itemize}[nosep]
    \item Core Question:
\end{itemize}

\textbf{In-Scope:}
\begin{itemize}[nosep]
    \item The Fundamental Question
    \item Core Theory: Strategy Design Model v1.0
    \item Axioms and Formal Definitions
    \item Key Results: Predictions from SDM-1.0
    \item Integration with EBF Framework
\end{itemize}

\textbf{Out-of-Scope (delegiert an andere Appendices):}
\begin{itemize}[nosep]
    \item HOW: Complementarities $\rightarrow$ Appendix UNMAPPED_HOW
    \item WHEN: Context $\rightarrow$ Appendix CTW
    \item Parameters $\rightarrow$ Appendix UNMAPPED_EST
    \item EEE Workflow used in SDM design $\rightarrow$ Appendix UNMAPPED_DSN
    \item HOW $\rightarrow$ Appendix UNMAPPED_HOW
\end{itemize}

\textbf{Constraints (Anwendungsgrenzen):}
\begin{itemize}[nosep]
    \item [TODO: Define when this appendix does NOT apply]
    \item [TODO: Define required prerequisites]
    \item [TODO: Define known limitations]
\end{itemize}

\textbf{Lieferobjekte (Deliverables):}
\begin{itemize}[nosep]
    \item 10 Table(s)
    \item 11 Equation(s)
    \item 3 Axiom(s)
    \item Worked Example(s)
\end{itemize}

\end{tcolorbox}


\section{The Fundamental Question}
\label{sec:sdm-fundamental}
%%%%%%%%%%%%%%%%%%%%%%%%%%%%%%%%%%%%%%%%%%%%%%%%%%%%%%%%%%%%%%%%%%%%%%%%%%%%%%%

\subsection{Why Corporate Strategy Design Matters}

Corporate strategy is arguably the most high-stakes organizational decision. A misdirected strategy can destroy decades of shareholder value, demoralize talented employees, and accelerate competitive decline. Yet strategy design remains largely an art rather than a science---guided by executive intuition, consulting frameworks, and historical analogies rather than evidence-based behavioral models.

The Evidence-Based Framework (EBF) offers a unique opportunity to formalize strategy design. Rather than treating strategy as a static document, SDM models strategy as a \textit{change journey}---a behavioral transformation unfolding across organizational units over multiple years. This process-oriented view reveals the true challenges: not articulating a good strategy, but \textit{implementing} one across thousands of employees with diverse incentives, capabilities, and beliefs.

\subsection{The Core Problem}

Strategic decisions exhibit three features that existing models fail to capture:

\begin{enumerate}
    \item \textbf{Complementarities:} Strategic dimensions interact. A financially attractive but strategically misaligned move (high F, low S) creates organizational friction and underexecution. The same move with high S creates synergistic buy-in. Standard additive utility functions miss this entirely.

    \item \textbf{Context Dependence:} The same strategy succeeds or fails depending on organizational context. A bold transformation strategy works during a growth cycle but fails during capital constraints. Strategy models must adapt to context ($\Psi$).

    \item \textbf{Dynamic Implementation:} Strategies unfold in phases. Organizations must move through Awareness → Readiness → Action → Embedding. Each phase has different bottlenecks (in Awareness, it's understanding; in Readiness, it's capability). Static models cannot capture this dynamism.
\end{enumerate}

\begin{tcolorbox}[colback=red!5!white, colframe=red!75!black,
    title=The Fundamental Question]
\textbf{How can we predict the success probability of a corporate strategy transformation, accounting for complementarity effects between strategic dimensions, context-dependent constraints, and multi-year implementation dynamics?}
\end{tcolorbox}

%%%%%%%%%%%%%%%%%%%%%%%%%%%%%%%%%%%%%%%%%%%%%%%%%%%%%%%%%%%%%%%%%%%%%%%%%%%%%%%
\section{Core Theory: Strategy Design Model v1.0}
\label{sec:sdm-theory}
%%%%%%%%%%%%%%%%%%%%%%%%%%%%%%%%%%%%%%%%%%%%%%%%%%%%%%%%%%%%%%%%%%%%%%%%%%%%%%%

\subsection{Strategic Context: Five Dimensions}
\label{subsec:sdm-context}

SDM incorporates five organizational context factors ($\Psi$) that modulate strategic decisions:

\begin{table}[htbp]
\centering
\caption{Context Factors ($\Psi$) in Strategy Design Model}
\label{tab:sdm-context}
\renewcommand{\arraystretch}{1.4}
\begin{tabular}{@{}lccp{5cm}@{}}
\toprule
\textbf{Factor} & \textbf{Symbol} & \textbf{Scale} & \textbf{Effect} \\
\midrule
Market Disruption & $\Psi_1$ & 0--1 & Higher $\Psi_1$ → More radical strategy needed \\
Capital Availability & $\Psi_2$ & 0--1 & Lower $\Psi_2$ → Conservative roadmap \\
Leadership Alignment & $\Psi_3$ & 0--1 & Lower $\Psi_3$ → Execution risk \\
Organizational Maturity & $\Psi_4$ & 0--1 & Lower $\Psi_4$ → Longer change journey \\
Regulatory Tightness & $\Psi_5$ & 0--1 & Higher $\Psi_5$ → More compliance overhead \\
\bottomrule
\end{tabular}
\end{table}

These factors are observable at strategy inception and can be estimated from organizational diagnostics. They directly modulate the strategic utility curve over the transformation timeline.

\subsection{Strategic Dimensions: FEPSDE for Corporate Strategy}
\label{subsec:sdm-dimensions}

SDM adapts the EBF FEPSDE framework (Appendix UNMAPPED_WAT) to five core strategic dimensions with organizational weights:

\begin{table}[htbp]
\centering
\caption{Strategic Dimensions and Weights in SDM}
\label{tab:sdm-dimensions}
\renewcommand{\arraystretch}{1.4}
\begin{tabular}{@{}lllp{5cm}@{}}
\toprule
\textbf{Dimension} & \textbf{Symbol} & \textbf{Weight $\beta_i$} & \textbf{Interpretation} \\
\midrule
Financial & $F$ & 0.25 & Profitability, ROI, cash generation \\
Environmental & $E$ & 0.15 & Sustainability, ESG, societal impact \\
Practical & $P$ & 0.20 & Operational feasibility, resource availability \\
Strategic & $S$ & 0.25 & Mission alignment, competitive advantage \\
Deliberative & $D$ & 0.10 & Analytical rigor, validation, scenario testing \\
\bottomrule
\end{tabular}
\end{table}

\textbf{Note:} Social dimension is implicit in Strategic and Environmental dimensions. For given strategy, scores range $C_i(t) \in [0,1]$ and evolve through BCJ phases.

\subsection{Complementarity Matrix: Seven $\gamma$ Parameters}
\label{subsec:sdm-gamma}

Strategic dimensions exhibit seven key complementarity effects ($\gamma_{ij}$), organized into positive synergies and negative trade-offs:

\begin{table}[htbp]
\centering
\caption{Complementarity Parameters ($\gamma$) in Strategy Design}
\label{tab:sdm-gamma}
\renewcommand{\arraystretch}{1.4}
\begin{tabular}{@{}llp{7cm}@{}}
\toprule
\textbf{Parameter} & \textbf{Value} & \textbf{Interpretation} \\
\midrule
\multicolumn{3}{l}{\textit{Positive Synergies (boost utility)}} \\
$\gamma_{FS}$ & +0.18 & Profit + Strategic alignment → market leadership \\
$\gamma_{SP}$ & +0.12 & Ambition + Execution → sustainable model \\
$\gamma_{SE}$ & +0.10 & Strategy + Impact → stronger org. commitment \\
$\gamma_{PD}$ & +0.10 & Feasibility + Rigor → validated implementation \\
$\gamma_{ED}$ & +0.08 & Sustainability + Analysis → durable advantage \\
\midrule
\multicolumn{3}{l}{\textit{Negative Trade-offs (reduce utility)}} \\
$\gamma_{FP}$ & -0.15 & High returns → higher complexity/risk \\
$\gamma_{FE}$ & -0.12 & Profit-max → environmental/social costs \\
\bottomrule
\end{tabular}
\end{table}

These parameters are derived from EBF appendix B (HOW: Complementarities) and organizational case evidence. They represent interaction effects that must be explicitly modeled for realistic strategy prediction.

\subsection{The BCJ (Behavioral Change Journey) Model}
\label{subsec:sdm-bcj}

Strategic transformations unfold across four distinct phases with characteristic organizational behaviors and bottlenecks:

\begin{table}[htbp]
\centering
\caption{The Four Phases of Strategic Change Journey}
\label{tab:sdm-phases}
\renewcommand{\arraystretch}{1.4}
\begin{tabular}{@{}llcll@{}}
\toprule
\textbf{Phase} & \textbf{Symbol} & \textbf{Timeline} & \textbf{Key Behavior} & \textbf{Bottleneck} \\
\midrule
Awareness & $\phi=0$ & 0--6 mo & Strategy articulated & Understanding \\
Readiness & $\phi=1/3$ & 6--18 mo & Org structures reset & Capability \\
Action & $\phi=2/3$ & 18--36 mo & Implementation in divisions & Commitment \\
Embedding & $\phi=1$ & 36--60 mo & New operating model locked & Sustainability \\
\bottomrule
\end{tabular}
\end{table}

Strategy success requires smooth progression through all four phases. Failure typically occurs when organizations:
\begin{itemize}
    \item Skip Readiness (implement too fast, capability gaps emerge)
    \item Stall in Action (early wins create false confidence, momentum lost)
    \item Fail Embedding (new model reverts to old behaviors under stress)
\end{itemize}

\subsection{Strategic Utility Function}
\label{subsec:sdm-utility}

The strategic utility at time $t$ (in years) is:

\begin{equation}
\label{eq:sdm-utility}
U_{\text{strategy}}(t) = \beta_0 + \sum_{i=1}^{5} \beta_i C_i(t, \phi(t)) + \sum_{i<j} \gamma_{ij} C_i(t) C_j(t) + \Psi_t \cdot \lambda
\end{equation}

where:
\begin{itemize}
    \item $\beta_0 = 0.50$ is the intercept (moderate baseline strategic support)
    \item $\beta_i$ are weights from Table \ref{tab:sdm-dimensions}
    \item $C_i(t, \phi(t))$ are strategic dimension scores evolving through phases
    \item $\gamma_{ij}$ are complementarity parameters from Table \ref{tab:sdm-gamma}
    \item $\Psi_t$ represents context modulation (typically $\Psi_t \in [-0.1, 0.2]$)
    \item $\lambda$ is the context sensitivity (typically $\lambda \approx 1.0$)
\end{itemize}

\textbf{Dimension Evolution:} Each dimension $C_i$ evolves as:

\begin{equation}
\label{eq:sdm-dimension-evolution}
C_i(t, \phi) = C_i^{\text{target}} \cdot f(\phi, \tau) + C_i^{\text{baseline}} \cdot (1 - f(\phi, \tau))
\end{equation}

where $f(\phi, \tau)$ is a sigmoid phase transition function:

\begin{equation}
f(\phi, \tau) = \frac{1}{1 + \exp(-\tau(\phi - 0.5))}
\end{equation}

Here $\tau$ is the organizational learning rate (speed of phase transition), typically $\tau \in [1, 3]$.

%%%%%%%%%%%%%%%%%%%%%%%%%%%%%%%%%%%%%%%%%%%%%%%%%%%%%%%%%%%%%%%%%%%%%%%%%%%%%%%
\section{Axioms and Formal Definitions}
\label{sec:sdm-axioms}
%%%%%%%%%%%%%%%%%%%%%%%%%%%%%%%%%%%%%%%%%%%%%%%%%%%%%%%%%%%%%%%%%%%%%%%%%%%%%%%

\begin{tcolorbox}[colback=gray!5!white, colframe=gray!75!black,
    title=Axiom SDM-1: Complementarity Necessity]

\textbf{Statement:} Strategic utility cannot be decomposed additively. For any two dimensions $C_i, C_j$ describing strategic moves:

\begin{equation}
U(C_i, C_j) \neq \alpha C_i + \beta C_j
\end{equation}

Rather, interactions exist with non-zero $\gamma_{ij}$.

\textbf{Interpretation:} Strategies are complex systems where dimensions interact synergistically or in conflict. A strategy high on Financial but low on Strategic alignment creates organizational friction beyond simple additive loss.

\textbf{Implication:} Strategy models must explicitly include complementarity terms. Frameworks that ignore $\gamma$ will systematically mispredic strategy success.

\textbf{Evidence:} Organizational case studies show high within-team disagreement when F and S are misaligned (Rumelt, 2011; Collis \& Rukstad, 2008).

\end{tcolorbox}

\begin{tcolorbox}[colback=gray!5!white, colframe=gray!75!black,
    title=Axiom SDM-2: Context Modulation]

\textbf{Statement:} The success probability of a given strategy is not fixed but depends on organizational context:

\begin{equation}
P(\text{Success} \mid \text{Strategy}) = f(C_1, C_2, C_3, C_4, C_5, \Psi_1, \Psi_2, \Psi_3, \Psi_4, \Psi_5)
\end{equation}

The same strategy has $P \approx 0.8$ during growth cycles but $P \approx 0.4$ during capital constraints.

\textbf{Interpretation:} Strategy design is not about absolute quality but about fit between strategy and context. Context factors must be estimated before strategy evaluation.

\textbf{Implication:} Strategy frameworks should include explicit context diagnostics (via $\Psi$ assessment) before strategy design begins.

\end{tcolorbox}

\begin{tcolorbox}[colback=gray!5!white, colframe=gray!75!black,
    title=Axiom SDM-3: Phase Progression]

\textbf{Statement:} Organizational strategies progress through discrete phases with distinct behaviors. Attempting to skip phases increases failure risk. Let $\phi(t) \in \{0, 1/3, 2/3, 1\}$ denote phase at time $t$:

\begin{equation}
\text{Skip Risk} = P(\text{Reversion} \mid \text{Skipped Phase}) \approx 0.6 \text{--} 0.8
\end{equation}

\textbf{Interpretation:} Change management literature consistently shows that rushing implementations without proper Readiness phase causes implementation failures. Organizations must invest in capability-building before Action phase.

\textbf{Evidence:} McKinsey Change Survey (2020) found 60-70% failure rate for rapid strategy implementations without adequate Readiness phase.

\end{tcolorbox}

%%%%%%%%%%%%%%%%%%%%%%%%%%%%%%%%%%%%%%%%%%%%%%%%%%%%%%%%%%%%%%%%%%%%%%%%%%%%%%%
\section{Key Results: Predictions from SDM-1.0}
\label{sec:sdm-results}
%%%%%%%%%%%%%%%%%%%%%%%%%%%%%%%%%%%%%%%%%%%%%%%%%%%%%%%%%%%%%%%%%%%%%%%%%%%%%%%

\subsection{Phase-Specific Success Probabilities}

By evaluating equation \ref{eq:sdm-utility} at each phase, we obtain:

\begin{table}[htbp]
\centering
\caption{Strategy Success Probability by Phase (Baseline Context)}
\label{tab:sdm-predictions}
\renewcommand{\arraystretch}{1.4}
\begin{tabular}{@{}lcccl@{}}
\toprule
\textbf{Phase} & \textbf{Timeline} & $\mathbf{U_{\text{strategy}}}$ & $\mathbf{P(\text{Success})}$ & \textbf{Recommendation} \\
\midrule
Awareness & 0--6 mo & 0.35--0.45 & 35--45\% & EVALUATE \\
Readiness & 6--18 mo & 0.45--0.60 & 45--60\% & PREPARE \\
Action & 18--36 mo & 0.60--0.75 & 60--75\% & EXECUTE \\
Embedding & 36--60 mo & 0.75--0.85 & 75--85\% & SUSTAIN \\
\bottomrule
\end{tabular}
\end{table}

\textbf{Interpretation:} These are baseline predictions assuming moderate context ($\Psi \approx 0.5$ for all factors). Context-adjusted predictions require specific organizational diagnostics.

\subsection{Context Sensitivity: Impact of Leadership Alignment}

Leadership alignment ($\Psi_3$) has particularly strong effect on Action phase success:

\begin{equation}
\Delta P(\text{Success}) = 0.15 \times (\Psi_3^{\text{high}} - \Psi_3^{\text{low}}) \approx 15\%
\end{equation}

Organizations with high leadership alignment ($\Psi_3=0.8$) achieve $P \approx 0.75$ in Action phase, while those with misaligned leadership ($\Psi_3=0.4$) achieve $P \approx 0.60$.

\subsection{Complementarity Impact: The $\gamma_{FS}$ Effect}

Strategies with strong Financial-Strategic alignment (high $F$, high $S$) show:

\begin{equation}
U_{\text{with } \gamma_{FS}} = 0.50 + 0.25F + 0.25S + 0.18 \cdot F \cdot S
\end{equation}

versus additive model:

\begin{equation}
U_{\text{additive}} = 0.50 + 0.25F + 0.25S
\end{equation}

For $F = S = 0.8$:
- With complementarity: $U = 0.50 + 0.20 + 0.20 + 0.115 = 0.915$
- Additive: $U = 0.50 + 0.20 + 0.20 = 0.90$

Complementarity adds $\approx 1.5\%$ marginal utility---small in isolation but critical across portfolio of strategic decisions.

%%%%%%%%%%%%%%%%%%%%%%%%%%%%%%%%%%%%%%%%%%%%%%%%%%%%%%%%%%%%%%%%%%%%%%%%%%%%%%%
\section{Integration with EBF Framework}
\label{sec:sdm-integration}
%%%%%%%%%%%%%%%%%%%%%%%%%%%%%%%%%%%%%%%%%%%%%%%%%%%%%%%%%%%%%%%%%%%%%%%%%%%%%%%

\subsection{Connection to Appendix UNMAPPED_HOW: Complementarities (HOW)}

SDM applies the general complementarity theory from Appendix UNMAPPED_HOW to the organizational strategy domain. The $\gamma$ matrix in Table \ref{tab:sdm-gamma} is an instance of the general EBF complementarity framework:

\begin{tcolorbox}[colback=yellow!5!white, colframe=yellow!75!black,
    title=Integration: SDM $\leftrightarrow$ Appendix UNMAPPED_HOW (HOW)]

\textbf{What SDM provides:}
\begin{itemize}[nosep]
\item Concrete instantiation of $\gamma$ parameters for strategic domain
\item Empirical estimates of interaction magnitudes (e.g., $\gamma_{FS} = +0.18$)
\item Evidence that complementarities are observable in organizational behavior
\end{itemize}

\textbf{What SDM requires from B:}
\begin{itemize}[nosep]
\item Theoretical justification for complementarity inclusion
\item Mathematical framework for $\gamma$ specification
\item Consistency requirements for $\gamma$ matrices (e.g., symmetry)
\end{itemize}

\textbf{Mathematical Interface:}

SDM's utility function (Eq. \ref{eq:sdm-utility}) is an application of EBF Appendix UNMAPPED_HOW's general form:

\begin{equation}
U(C) = \beta_0 + \sum_i \beta_i C_i + \underbrace{\sum_{i<j} \gamma_{ij} C_i C_j}_{\text{Appendix UNMAPPED_HOW}} + \lambda \Psi
\end{equation}

\end{tcolorbox}

\subsection{Connection to Appendix CTW: Context (WHEN)}

SDM operationalizes the general context framework from Appendix CTW specifically for organizational strategy. The five context factors ($\Psi_1$ through $\Psi_5$) are domain-specific instantiations of the general EBF context concept:

\begin{tcolorbox}[colback=yellow!5!white, colframe=yellow!75!black,
    title=Integration: SDM $\leftrightarrow$ Appendix CTW (WHEN)]

\textbf{What SDM provides:}
\begin{itemize}[nosep]
\item Measurement protocol for $\Psi$ factors in organizational context
\item Quantitative modulation functions showing how $\Psi$ affects $U$
\item Evidence of context-dependent success probabilities
\end{itemize}

\textbf{What SDM requires from V:}
\begin{itemize}[nosep]
\item Theoretical framework for why context matters
\item Axioms governing context modulation
\item Consistency across domains
\end{itemize}

\end{tcolorbox}

\subsection{Connection to Appendix UNMAPPED_WAT: Strategic Dimensions (WHAT)}

SDM adapts the FEPSDE framework from Appendix UNMAPPED_WAT to the organizational strategy domain. The five strategic dimensions in Table \ref{tab:sdm-dimensions} map to EBF's core utility dimensions:

\begin{table}[htbp]
\centering
\caption{Mapping: SDM Dimensions $\leftrightarrow$ EBF FEPSDE}
\label{tab:sdm-fepsde-mapping}
\renewcommand{\arraystretch}{1.3}
\begin{tabular}{@{}lll@{}}
\toprule
\textbf{SDM Dimension} & \textbf{EBF FEPSDE} & \textbf{Note} \\
\midrule
Financial & F & Direct correspondence \\
Environmental & E & Direct correspondence \\
Practical & P & Direct correspondence \\
Strategic & S & Combines strategic + social preferences \\
Deliberative & D & Direct correspondence \\
\bottomrule
\end{tabular}
\end{table}

%%%%%%%%%%%%%%%%%%%%%%%%%%%%%%%%%%%%%%%%%%%%%%%%%%%%%%%%%%%%%%%%%%%%%%%%%%%%%%%
% PART 3: APPLICATION
%%%%%%%%%%%%%%%%%%%%%%%%%%%%%%%%%%%%%%%%%%%%%%%%%%%%%%%%%%%%%%%%%%%%%%%%%%%%%%%

%%%%%%%%%%%%%%%%%%%%%%%%%%%%%%%%%%%%%%%%%%%%%%%%%%%%%%%%%%%%%%%%%%%%%%%%%%%%%%%
\section{Worked Example: Digital Transformation Strategy}
\label{sec:sdm-example}
%%%%%%%%%%%%%%%%%%%%%%%%%%%%%%%%%%%%%%%%%%%%%%%%%%%%%%%%%%%%%%%%%%%%%%%%%%%%%%%

\subsection{Case: Large Bank's Digital Transformation (2024--2029)}

\paragraph{Setup.}

A global financial services institution with 50,000 employees must transform its core business from branch-based retail banking to digital-first operations. The board has committed \$2 billion over 5 years to this transformation. Key question: What is the predicted success probability, and how does leadership alignment affect outcomes?

\paragraph{Context Assessment ($\Psi$ Factors).}

\begin{table}[htbp]
\centering
\caption{Context Assessment: Bank Digital Transformation}
\label{tab:sdm-case-context}
\renewcommand{\arraystretch}{1.3}
\begin{tabular}{@{}lll@{}}
\toprule
\textbf{Context Factor} & \textbf{Assessment} & \textbf{Score} \\
\midrule
Market Disruption ($\Psi_1$) & High tech competition (Fintechs, BigTech) & 0.70 \\
Capital Availability ($\Psi_2$) & Well-capitalized, post-stress test & 0.75 \\
Leadership Alignment ($\Psi_3$) & Mixed: CEO aligned, CFO skeptical & 0.55 \\
Organizational Maturity ($\Psi_4$) & Legacy IT culture, limited tech talent & 0.40 \\
Regulatory Tightness ($\Psi_5$) & Stringent data privacy, compliance & 0.65 \\
\bottomrule
\end{tabular}
\end{table}

\paragraph{Strategic Dimension Assessment.}

Based on strategy document analysis and leadership interviews, dimension scores are:

\begin{table}[htbp]
\centering
\caption{Strategic Dimensions: Bank Digital Transformation}
\label{tab:sdm-case-dimensions}
\renewcommand{\arraystretch}{1.3}
\begin{tabular}{@{}lll@{}}
\toprule
\textbf{Dimension} & \textbf{Assessment} & \textbf{Score} \\
\midrule
Financial ($F$) & Strong ROI case, but 3-year investment period & 0.75 \\
Environmental ($E$) & Green finance opportunities & 0.60 \\
Practical ($P$) & Challenging: requires tech hiring, org restructure & 0.50 \\
Strategic ($S$) & Clear alignment with market necessity & 0.80 \\
Deliberative ($D$) & Extensive scenario testing, validated roadmap & 0.75 \\
\bottomrule
\end{tabular}
\end{table}

\paragraph{Utility Calculation by Phase.}

Using Equation \ref{eq:sdm-utility} with $\lambda = 1.0$:

\textbf{Phase 1: Awareness (Month 0--6):}

\begin{align}
U_{\text{Awareness}} &= 0.50 + (0.25 \times 0.75 + 0.15 \times 0.60 + 0.20 \times 0.50 + 0.25 \times 0.80 + 0.10 \times 0.75) \\
&\quad + (0.18 \times 0.75 \times 0.80 + 0.12 \times 0.80 \times 0.50 + \ldots) \\
&\quad + (0.55 - 0.50) \\
&= 0.50 + 0.66 + 0.13 + 0.05 = 1.34
\end{align}

Clipped to $[0, 1]$ for probability: $P(\text{Awareness}) = 0.95$.

\textbf{Phase 2: Readiness (Month 6--18):}

Dimension scores evolve; organizational maturity becomes limiting factor:

\begin{align}
U_{\text{Readiness}} &= 0.50 + 0.62 + 0.09 + 0.02 = 0.83 \\
P(\text{Readiness}) &\approx 0.83
\end{align}

The reduced $P$ reflects maturity gaps in Readiness phase (capability building more difficult than strategy articulation).

\textbf{Phase 3: Action (Month 18--36):}

Implementation across retail divisions; leadership alignment critical:

With $\Psi_3 = 0.55$ (moderate misalignment):
\begin{align}
U_{\text{Action}} &= 0.50 + 0.58 + 0.08 - 0.02 = 0.64 \\
P(\text{Action}) &\approx 0.68
\end{align}

\textbf{Scenario: Improved Leadership Alignment ($\Psi_3 = 0.75$):}

If board resolves CFO concerns and achieves alignment:
\begin{align}
U_{\text{Action}} &= 0.64 + (0.75 - 0.55) \approx 0.84 \\
P(\text{Action}) &\approx 0.84 \quad (\text{+16\% improvement})
\end{align}

\paragraph{Result.}

Baseline transformation success probability: $P \approx 0.68$ (conditional on full 4-phase completion). Main risks: organizational maturity gaps ($\Psi_4 = 0.40$) and leadership misalignment ($\Psi_3 = 0.55$).

\paragraph{Discussion.}

The model identifies two high-leverage interventions:

\begin{enumerate}
    \item \textbf{Increase Leadership Alignment ($\Psi_3$):} Address CFO concerns with pilot evidence. Potential +16\% success boost.

    \item \textbf{Accelerate Organizational Maturity ($\Psi_4$):} Hire external tech talent, establish tech COE. Longer-term but essential for Action phase.
\end{enumerate}

Without these interventions, transformation success remains below 70\%. With them, success probability exceeds 80\%.

%%%%%%%%%%%%%%%%%%%%%%%%%%%%%%%%%%%%%%%%%%%%%%%%%%%%%%%%%%%%%%%%%%%%%%%%%%%%%%%
\section{Implications for Practice}
\label{sec:sdm-practice}
%%%%%%%%%%%%%%%%%%%%%%%%%%%%%%%%%%%%%%%%%%%%%%%%%%%%%%%%%%%%%%%%%%%%%%%%%%%%%%%

\begin{tcolorbox}[colback=green!5!white, colframe=green!75!black,
    title=Practical Implications for Strategy Design]

\begin{enumerate}

\item \textbf{Diagnose Context Before Strategy Design:} Assess all five $\Psi$ factors (market disruption, capital, leadership alignment, org maturity, regulatory tightness) before committing to strategy. Context constrains feasibility more than strategy quality.

\item \textbf{Demand Complementarity Alignment:} Strategies must score high on \textit{both} Financial ($F$) and Strategic ($S$) dimensions (or other key pairs). Misaligned strategies generate internal friction and underexecution. Use $\gamma$ matrix to identify misalignments early.

\item \textbf{Invest in Readiness Phase:} The biggest mistake is rushing from Awareness to Action. Allocate 6--18 months and 20--30\% of transformation budget to Readiness (capability building, org structure, systems). This phase determines Action success.

\item \textbf{Use BCJ Model for Phased Rollout:} Rather than all-at-once implementation, sequence strategy across four phases. Each phase has different success drivers. In Awareness, focus on understanding; in Readiness, on capability; in Action, on execution; in Embedding, on sustainability.

\item \textbf{Monitor Context Continuously:} $\Psi$ factors change. If capital becomes constrained during implementation, adjust strategy scope. If leadership alignment deteriorates, invest in realignment before Action phase. Treat strategy as adaptive, not fixed.

\end{enumerate}

\end{tcolorbox}

%%%%%%%%%%%%%%%%%%%%%%%%%%%%%%%%%%%%%%%%%%%%%%%%%%%%%%%%%%%%%%%%%%%%%%%%%%%%%%%
% PART 4: BACK MATTER
%%%%%%%%%%%%%%%%%%%%%%%%%%%%%%%%%%%%%%%%%%%%%%%%%%%%%%%%%%%%%%%%%%%%%%%%%%%%%%%

%%%%%%%%%%%%%%%%%%%%%%%%%%%%%%%%%%%%%%%%%%%%%%%%%%%%%%%%%%%%%%%%%%%%%%%%%%%%%%%
\section{Summary}
\label{sec:sdm-summary}
%%%%%%%%%%%%%%%%%%%%%%%%%%%%%%%%%%%%%%%%%%%%%%%%%%%%%%%%%%%%%%%%%%%%%%%%%%%%%%%

\begin{tcolorbox}[colback=green!10!white, colframe=green!75!black,
    title=Appendix UNMAPPED_COR: Strategy Design Model v1.0 --- Summary]

\textbf{Core Question:} How do organizations successfully transition from current to target strategy over 3--5 years?

\textbf{Main Contribution:}
\begin{itemize}[nosep]
    \item Process-based model treating strategy as 4-phase behavioral change journey
    \item Explicit complementarity effects between strategic dimensions ($\gamma$ matrix)
    \item Context-dependent success predictions ($\Psi$ factors)
    \item Practical framework for strategy diagnostics and risk assessment
\end{itemize}

\textbf{Key Outputs:}
\begin{itemize}[nosep]
    \item \textbf{Utility Function ($U$):} Predicts strategic success probability given context and dimensions
    \item \textbf{Complementarity Matrix ($\gamma$):} Quantifies synergies and trade-offs between dimensions
    \item \textbf{Context Modulation ($\Psi$):} Adjusts predictions based on organizational environment
    \item \textbf{Phase Roadmap (BCJ):} Sequences strategy implementation with phase-specific bottlenecks
\end{itemize}

\textbf{Integration:} SDM operationalizes EBF theory (complementarities from B, context from V, dimensions from C) in the organizational strategy domain, connecting abstract framework to concrete strategic decisions.

\end{tcolorbox}

%%%%%%%%%%%%%%%%%%%%%%%%%%%%%%%%%%%%%%%%%%%%%%%%%%%%%%%%%%%%%%%%%%%%%%%%%%%%%%%
\section{Glossary of Symbols}
\label{sec:sdm-glossary}
%%%%%%%%%%%%%%%%%%%%%%%%%%%%%%%%%%%%%%%%%%%%%%%%%%%%%%%%%%%%%%%%%%%%%%%%%%%%%%%

\begin{tcolorbox}[colback=blue!5!white, colframe=blue!75!black]
\textbf{Central Reference:} For complete symbol definitions, see
\textbf{Appendix UNMAPPED_GLS} (\textit{Glossary and Notation}, \S\ref{app:glossary}).

This section lists symbols introduced or specialized in SDM.
\end{tcolorbox}

\vspace{0.5em}

\begin{table}[htbp]
\centering
\caption{Symbols Introduced in Appendix UNMAPPED_COR (SDM)}
\label{tab:sdm-symbols}
\renewcommand{\arraystretch}{1.3}
\begin{tabular}{@{}clp{5.5cm}c@{}}
\toprule
\textbf{Symbol} & \textbf{Name} & \textbf{Definition} & \textbf{In G?} \\
\midrule
$U_{\text{strategy}}(t)$ & Strategic Utility & Probability of strategy success at time $t$ & NEW \\
$C_i$ & Dimension Score & Score for dimension $i$ (Financial, Environmental, etc.) & \checkmark \\
$\gamma_{ij}$ & Complementarity & Interaction effect between dimensions $i$ and $j$ & \checkmark \\
$\Psi_k$ & Context Factor & $k$-th organizational context factor & \checkmark \\
$\phi(t)$ & Phase Progress & Position in BCJ (0=Awareness, 1=Embedding) & NEW \\
$\tau$ & Learning Rate & Speed of organizational phase progression & NEW \\
$\beta_0$ & Intercept & Baseline strategic support (0.50) & NEW \\
\bottomrule
\end{tabular}
\end{table}

%%%%%%%%%%%%%%%%%%%%%%%%%%%%%%%%%%%%%%%%%%%%%%%%%%%%%%%%%%%%%%%%%%%%%%%%%%%%%%%
\section{Critical Foundations and Objections}
\label{sec:sdm-critical}
%%%%%%%%%%%%%%%%%%%%%%%%%%%%%%%%%%%%%%%%%%%%%%%%%%%%%%%%%%%%%%%%%%%%%%%%%%%%%%%

\subsection{Objection 1: ``Strategy is Art, Not Science''}

\begin{tcolorbox}[colback=red!5!white, colframe=red!50!black,
    title=Objection: Strategy Cannot Be Formalized]
Traditional strategy scholars argue strategy formulation is fundamentally creative and context-specific, resisting quantification.
\end{tcolorbox}

\textbf{Response:} SDM does not claim strategy formulation is mechanical. Rather, it formalizes the \textit{evaluation} and \textit{implementation} of strategies once articulated. The creative work (identifying opportunities, framing alternatives) remains. But after creation, SDM provides a quantitative framework for:
\begin{itemize}
    \item Predicting implementation success given organizational context
    \item Identifying complementarity misalignments early
    \item Sequencing implementation phases optimally
\end{itemize}

This is analogous to architecture: designing a building is creative, but structural engineering then quantifies whether it will stand.

\subsection{Objection 2: ``Parameters Are Arbitrary''}

\begin{tcolorbox}[colback=red!5!white, colframe=red!50!black,
    title=Objection: Complementarity Values ($\gamma$) Are Ad-hoc]
Critics argue that $\gamma$ parameter values (e.g., $\gamma_{FS} = +0.18$) lack empirical validation.
\end{tcolorbox}

\textbf{Response:} Fair criticism. SDM-1.0 uses conservative estimates based on:
\begin{itemize}
    \item EBF Appendix UNMAPPED_EST (Parameter Repository) literature review
    \item Organizational case studies (strategic initiative success/failure analysis)
    \item Limited organizational dataset ($n \approx 40$ strategic initiatives)
\end{itemize}

These parameters should be treated as Bayesian priors, not fixed values. Phase 1.1 roadmap includes parameter calibration with larger empirical dataset. Organizations implementing SDM should adjust $\gamma$ based on their industry and history.

\subsection{Objection 3: ``This Ignores Leadership and Politics''}

\begin{tcolorbox}[colback=red!5!white, colframe=red!50!black,
    title=Objection: Model Omits Political Dynamics]
Organizational politics and leadership personalities often trump formal strategy frameworks.
\end{tcolorbox}

\textbf{Response:} Acknowledged. SDM v1.0 models organizational maturity ($\Psi_4$) and leadership alignment ($\Psi_3$) as scalar factors, abstracting away political detail. This is a conscious simplification for analytical tractability.

More sophisticated models (Phase 2.0 roadmap) will incorporate:
\begin{itemize}
    \item Behavioral biases (organizational herding, loss aversion)
    \item Political coalitions and veto points
    \item Executive compensation alignment with strategy
\end{itemize}

For now, SDM acknowledges politics indirectly through $\Psi_3$ (leadership alignment) and provides guidance on improving it through the interventions box.

%%%%%%%%%%%%%%%%%%%%%%%%%%%%%%%%%%%%%%%%%%%%%%%%%%%%%%%%%%%%%%%%%%%%%%%%%%%%%%%
\section{Open Issues and Research Agenda}
\label{sec:sdm-open}
%%%%%%%%%%%%%%%%%%%%%%%%%%%%%%%%%%%%%%%%%%%%%%%%%%%%%%%%%%%%%%%%%%%%%%%%%%%%%%%

\begin{table}[htbp]
\centering
\caption{Research Roadmap for Strategy Design Model}
\label{tab:sdm-roadmap}
\renewcommand{\arraystretch}{1.3}
\begin{tabular}{@{}clccl@{}}
\toprule
\textbf{\#} & \textbf{Issue} & \textbf{Priority} & \textbf{Status} & \textbf{Requirements} \\
\midrule
1 & Empirical validation of $\gamma$ parameters & 🔴 HIGH & Pending & 100+ strategic initiatives \\
2 & Cross-industry generalization & 🔴 HIGH & Pending & 5+ industry datasets \\
3 & Integration of behavioral biases & 🟡 MEDIUM & Planned & Q2 2026 \\
4 & Real-time context monitoring & 🟡 MEDIUM & Planned & Q2 2026 \\
5 & Political coalition modeling & 🟡 MEDIUM & Planned & Q3 2026 \\
6 & AI-assisted strategy formulation & 🟢 LOW & Future & 2027+ \\
7 & Multi-stakeholder utility aggregation & 🟢 LOW & Future & Theoretical work \\
\bottomrule
\end{tabular}
\end{table}

%%%%%%%%%%%%%%%%%%%%%%%%%%%%%%%%%%%%%%%%%%%%%%%%%%%%%%%%%%%%%%%%%%%%%%%%%%%%%%%
\section{References}
\label{sec:sdm-references}
%%%%%%%%%%%%%%%%%%%%%%%%%%%%%%%%%%%%%%%%%%%%%%%%%%%%%%%%%%%%%%%%%%%%%%%%%%%%%%%

Key references for Strategy Design Model:

\nocite{bcm_master}

\begin{thebibliography}{99}

\bibitem{collis2008} Collis, D. J., \& Rukstad, M. E. (2008). Can You Say What Your Strategy Is? \textit{Harvard Business Review}, 86(4), 82--90.

\bibitem{rumelt2011} Rumelt, R. P. (2011). \textit{Good Strategy / Bad Strategy: The Difference and Why It Matters}. Crown Business.

\bibitem{mckinsey2020} McKinsey \& Company (2020). Winning Through Transformation: A Change Readiness Guide. McKinsey Survey on Organizational Capability.

\bibitem{andrews1971} Andrews, K. R. (1971). \textit{The Concept of Corporate Strategy}. Homewood, IL: Dow Jones-Irwin.

\bibitem{porter1980} Porter, M. E. (1980). \textit{Competitive Strategy}. Free Press.

\end{thebibliography}

%%%%%%%%%%%%%%%%%%%%%%%%%%%%%%%%%%%%%%%%%%%%%%%%%%%%%%%%%%%%%%%%%%%%%%%%%%%%%%%
% CROSS-REFERENCE FOOTER
%%%%%%%%%%%%%%%%%%%%%%%%%%%%%%%%%%%%%%%%%%%%%%%%%%%%%%%%%%%%%%%%%%%%%%%%%%%%%%%

\vspace{1em}
\hrule
\vspace{0.5em}

\noindent\textit{Cross-references:}
Appendix UNMAPPED_HOW (HOW: Complementarities),
Appendix UNMAPPED_WAT (WHAT: Utility Dimensions),
Appendix CTW (WHEN: Context),
Appendix UNMAPPED_EST (Parameter Repository),
Appendix UNMAPPED_DSN (EEE Workflow).

\noindent\textit{Master Bibliography:} \texttt{bcm\_master\_references.bib}

\noindent\textit{Related Domain Appendices:} Appendix UNMAPPED_BNK (SBDMS: Bank Strategic Decisions), Appendix MUL1 (Personal Financial Planning)

% =============================================================================
% END OF APPENDIX BI
% =============================================================================
